\section{Cultura de equipo y capacitación}
\subsection{Estrategia de implementación cultural}
Graham enfatiza: cuando el Agent no tiene un desempeño sólido, las personas dudan de su valor; por eso hay que difundir dentro del equipo la cultura de «Agent + humano», dando como ejemplo que el Agent se encarga de tareas repetitivas, mientras el humano conserva la revisión final. En cada retrospectiva de sprint hay que mostrar los resultados del Agent.

\begin{enumerate}
    \item Compartir semanalmente casos de éxito del Agent (por ejemplo, el Agent genera un diff y pasa 3 pruebas).
    \item Elaborar un documento de «Agent Etiquette» que precise cuándo se puede invocar al Agent y en qué casos es obligatoria la revisión humana.
    \item Establecer un programa de Mentor, en el que un ingeniero con experiencia guíe a los nuevos integrantes antes de que se familiaricen con las herramientas del Agent.
\end{enumerate}

\subsection{Capacitación y guías}
El material de capacitación debe incluir terminología básica (Plan-Act-Verify/Observability Stream), los scripts más utilizados y advertencias de riesgo. Al terminar la capacitación se puede organizar una práctica real: hacer que los alumnos, en un sandbox, dejen que el Agent procese un Issue, mientras observan a un lado y registran los puntos de intervención.

\begin{warningbox}{Retos comunes de la capacitación}
Los nuevos integrantes tienden a depender demasiado del Agent; se recomienda establecer, en las prácticas reales, un umbral de «revisión humana» para asegurarse de que sepan verificar la salida del Agent.
\end{warningbox}

\begin{knowledgebox}{Transmisión de conocimiento}
Registrar en documentos y videos el proceso de cada colaboración entre humanos y máquinas, para facilitar que los equipos compartan la experiencia entre sí y reducir el riesgo de que «solo una persona sepa usar el Agent».
\end{knowledgebox}

\subsection{Resumen de la sección}
\begin{itemize}
    \item La cultura, la capacitación y las prácticas reales sostienen juntas la implementación a largo plazo del Agent.
    \item Una etiquette clara y un mecanismo de retrospectiva pueden eliminar la desconfianza hacia el Agent.
\end{itemize}

